\documentclass[12pt,a4paper]{article}

% =========================
% Pacotes essenciais
% =========================
\usepackage[a4paper,margin=2.5cm]{geometry}
\usepackage[T1]{fontenc}
\usepackage[utf8]{inputenc}
\usepackage[brazil]{babel}
\usepackage{lmodern}
\usepackage{setspace}
\usepackage{microtype}
\usepackage{hyperref}
\usepackage{enumitem}
\usepackage{fancyhdr}

\hypersetup{
  colorlinks=true,
  linkcolor=black,
  urlcolor=blue,
  citecolor=black
}

\setstretch{1.15}
\setlength{\parindent}{0pt}
\setlength{\parskip}{6pt}

\pagestyle{fancy}
\fancyhf{}
\lhead{Manual do Professor — CV721A}
\rhead{\thepage}
\cfoot{}

% =========================
% Metadados do documento
% =========================
\newcommand{\Disciplina}{Fundações (CV721A)}
\newcommand{\Curso}{Engenharia Civil — UNICAMP}
\newcommand{\VersaoManual}{v1.0}
\newcommand{\DataManual}{\today}

\title{\textbf{Manual do Professor}\\
Sistema de Apoio ao Trabalho Docente\\
\Disciplina}
\author{}
\date{\DataManual\ (\VersaoManual)}

\begin{document}
\maketitle

\section{Apresentação}
Este manual orienta o uso do \textbf{Sistema de Apoio ao Trabalho Docente}, desenvolvido para auxiliar o professor no gerenciamento de:
\begin{itemize}[noitemsep]
  \item presença diária dos alunos;
  \item consolidação de notas (provas, trabalhos e atividades interativas);
  \item exportação de dados acadêmicos para uso no sistema institucional da UNICAMP.
\end{itemize}

O sistema foi projetado para \textbf{reduzir tarefas operacionais}, mantendo o foco do professor na condução da aula.

\textbf{Observação:} o sistema \textbf{não substitui} o Google Classroom nem sistemas institucionais da UNICAMP. Ele atua como uma \textbf{ferramenta complementar}, dedicada ao trabalho docente em sala.

\section{Visão Geral do Sistema}
O sistema é acessado via navegador (computador ou tablet do professor) e organiza as principais tarefas docentes em um \textbf{painel único}, composto pelos seguintes módulos:
\begin{itemize}[noitemsep]
  \item Aula de Hoje;
  \item Frequência;
  \item Notas \& Avaliações;
  \item Importar Kahoot;
  \item Exportar UNICAMP.
\end{itemize}

\section{Requisitos para Uso}
\subsection*{Para o professor}
\begin{itemize}[noitemsep]
  \item notebook ou computador com acesso à internet;
  \item navegador web atualizado (Chrome, Firefox, Edge ou Safari);
  \item projetor ou tela para exibir o QR Code.
\end{itemize}

\subsection*{Para os alunos}
\begin{itemize}[noitemsep]
  \item smartphone com câmera;
  \item acesso à internet (Wi-Fi ou dados móveis).
\end{itemize}

\section{Módulo ``Aula de Hoje''}
Este é o \textbf{módulo principal} do sistema, utilizado em todas as aulas.

\subsection{Material da Aula (opcional)}
O professor pode, se desejar:
\begin{itemize}[noitemsep]
  \item inserir um \textbf{link para o material da aula} (Google Drive, OneDrive etc.);
  \item ou fazer upload de um arquivo (\texttt{.pptx} ou \texttt{.pdf}).
\end{itemize}
O botão \textbf{``Abrir aula''} abre o material em nova aba, funcionando como atalho organizado para o PowerPoint. Este recurso é opcional e não interfere no controle de presença.

\subsection{Gerar Presença por QR Code}
\begin{enumerate}[label=\arabic*)]
  \item No início da aula, clique em \textbf{``Gerar QR de Presença''}.
  \item Um QR Code será exibido na tela.
  \item Projete o QR Code para os alunos.
  \item O QR Code possui validade temporária (ex.: 3--5 minutos).
\end{enumerate}
Durante o período de validade, os alunos acessam o QR, informam \textbf{RA} e \textbf{Nome}, e a presença é registrada automaticamente.

\subsection{Acompanhamento em tempo real}
Enquanto o QR estiver ativo, o professor pode acompanhar:
\begin{itemize}[noitemsep]
  \item número de alunos presentes;
  \item lista de presenças registradas (RA, Nome e Horário).
\end{itemize}
Esse acompanhamento é opcional e não interfere no andamento da aula.

\subsection{Encerrar Aula}
Ao final do período de presença:
\begin{enumerate}[label=\arabic*)]
  \item Clique em \textbf{``Encerrar Aula''}.
  \item O QR Code é invalidado.
  \item Nenhum novo check-in é aceito.
\end{enumerate}
Isso evita registros fora do horário da aula.

\section{Módulo ``Frequência''}
Este módulo permite acompanhar a frequência ao longo do semestre, incluindo:
\begin{itemize}[noitemsep]
  \item presença por aluno;
  \item presença por data/aula;
  \item cálculo de faltas;
  \item correções pontuais (exceções) quando necessário.
\end{itemize}

\section{Módulo ``Notas \& Avaliações''}
Este módulo centraliza todas as notas da disciplina.

\subsection{Criar avaliações}
O professor pode criar avaliações do tipo: Prova, Trabalho, Kahoot e outras atividades. Para cada avaliação, define-se: nome, tipo e peso.

\subsection{Inserção de notas}
As notas podem ser inseridas manualmente ou importadas (ver módulo Kahoot). As médias parciais podem ser calculadas automaticamente conforme a configuração adotada.

\section{Módulo ``Importar Kahoot''}
Este módulo simplifica o uso de atividades interativas como nota.
\begin{enumerate}[label=\arabic*)]
  \item Aplicar a atividade no Kahoot.
  \item Exportar o resultado em formato \textbf{CSV}.
  \item Acessar o módulo \textbf{Importar Kahoot}.
  \item Fazer upload do arquivo CSV.
  \item Associar o arquivo a uma avaliação existente ou criar uma nova.
  \item Conferir a pré-visualização dos dados.
  \item Confirmar a importação.
\end{enumerate}
O sistema cruza dados preferencialmente por \textbf{RA} e registra as notas automaticamente. Caso algum aluno não seja reconhecido, o sistema sinaliza e permite associação manual.

\section{Módulo ``Exportar UNICAMP''}
Ao final do semestre, o professor pode exportar \textbf{frequência} e \textbf{notas} em \textbf{CSV} para facilitar a transferência ao sistema institucional da UNICAMP, reduzindo tempo e erros manuais.

\section{Boas Práticas de Uso}
\begin{itemize}[noitemsep]
  \item Gerar o QR Code sempre no início da aula.
  \item Utilizar tempo curto de validade para evitar registros indevidos.
  \item Encerrar a aula após o período de check-in.
  \item Importar resultados do Kahoot logo após a atividade.
  \item Padronizar RA como identificador principal.
\end{itemize}

\section{Considerações Finais}
O sistema foi desenvolvido com foco em simplicidade, confiabilidade e integração com práticas já consolidadas. Sua evolução deve ser guiada pelas necessidades reais do professor, podendo incorporar novas funcionalidades ao longo do semestre.

\section*{Contato e Suporte}
Em caso de dúvidas, sugestões ou problemas, registrar a ocorrência e encaminhar para a equipe responsável pelo projeto piloto.

\end{document}
